
\documentclass{article}
\usepackage{colortbl}
\usepackage{makecell}
\usepackage{multirow}
\usepackage{supertabular}

\begin{document}

\newcounter{utterance}

\centering \large Interaction Transcript for game `codenames', experiment `risk\_high', episode 2 with claude{-}opus{-}4{-}8{-}azure.
\vspace{24pt}

{ \footnotesize  \setcounter{utterance}{1}
\setlength{\tabcolsep}{0pt}
\begin{supertabular}{c@{$\;$}|p{.15\linewidth}@{}p{.15\linewidth}p{.15\linewidth}p{.15\linewidth}p{.15\linewidth}p{.15\linewidth}}
    \# & $\;$A & \multicolumn{4}{c}{Game Master} & $\;\:$B\\
    \hline

    \theutterance \stepcounter{utterance}  
    & & & \multicolumn{2}{p{0.3\linewidth}}{
        \cellcolor[rgb]{0.9,0.9,0.9}{
            \makecell[{{p{\linewidth}}}]{
                \texttt{\tiny{[GM$|$GM]}}
                {'hidden': {'team': ['plastic', 'leu', 'geniu', 'complot', 'spaţiu', 'rădăcină', 'canada', 'cui', 'şoarece'], 'innocent': ['bloc', 'moarte', 'pitic'], 'opponent': ['caracatiţă', 'germania', 'vânt', 'stadion', 'prinţesă', 'pinguin', 'mamut', 'card'], 'assassin': ['bivol', 'împuşcat', 'rază', 'berlin', 'model']}, 'revealed': {'team': {'team': [], 'innocent': [], 'opponent': [], 'assassin': []}, 'opponent': {'team': [], 'innocent': [], 'opponent': [], 'assassin': []}}}
            }
        }
    }
    & & \\ \\

    \theutterance \stepcounter{utterance}  
    & & \multicolumn{4}{p{0.6\linewidth}}{
        \cellcolor[rgb]{0.9,0.9,0.9}{
            \makecell[{{p{\linewidth}}}]{
                \texttt{\tiny{[P1$\langle$GM]}}
                \texttt{Furnizați un indiciu dintr{-}un singur cuvânt care se referă la unele dintre cuvintele echipei dvs. din lista de mai jos. Puteți alege la câte și la ce cuvinte se referă indiciul dumneavoastră.} \\
\texttt{Scopul este ca partenerul tău să ghicească toate cuvintele echipei tale pe baza indiciilor pe care i le dai, înainte ca cealaltă echipă să le descopere pe toate. Dacă partenerul dvs. ghicește un cuvânt "adversar" sau un cuvânt "distractor", rândul său de ghicit se încheie, iar echipa adversă are ocazia să facă un tur. Dacă partenerul tău ghicește un cuvânt "ucigaș", jocul se încheie și echipa ta pierde imediat.} \\
\texttt{Cu cât indiciul tău se referă la mai multe cuvinte, cu atât câștigi mai repede jocul, dar cu atât asocierile pot fi mai greu de ghicit pentru partenerul tău, așa că alege{-}ți cu înțelepciune cuvântul indiciu. Cuvântul indiciu trebuie să fie legat semantic de cuvintele țintă, nu poate fi unul dintre cuvintele din liste sau să conțină părți ale acestora.} \\
\\ 
\texttt{Furnizați întotdeauna indiciul cu un singur cuvânt și lista de cuvinte țintă, separate prin virgulă, în următorul format și faceți răspunsurile cât mai scurte posibil; nu includeți niciodată alt text decât este necesar în acest formular:} \\
\\ 
\texttt{INDICIU: <CUVÂNT>} \\
\texttt{ȚINTE: <ȚINTE>} \\
\\ 
\texttt{Cuvintele echipei tale sunt: plastic, leu, geniu, complot, spaţiu, rădăcină, canada, cui, şoarece.} \\
\texttt{Cuvintele echipei adversarului dvs. sunt: caracatiţă, germania, vânt, stadion, prinţesă, pinguin, mamut, card.} \\
\texttt{Cuvintele de distragere a atenției sunt: bloc, moarte, pitic.} \\
\texttt{Cuvintele de ucidere sunt: bivol, împuşcat, rază, berlin, model.} \\
            }
        }
    }
    & \\ \\

    \theutterance \stepcounter{utterance}  
    & \multicolumn{4}{p{0.6\linewidth}}{
        \cellcolor[rgb]{0.9,0.9,0.9}{
            \makecell[{{p{\linewidth}}}]{
                \texttt{\tiny{[P1$\rangle$GM]}}
                \texttt{INDICIU: calculator} \\
\texttt{ȚINTE: şoarece, spaţiu} \\
            }
        }
    }
    & & \\ \\

    \theutterance \stepcounter{utterance}  
    & & & \multicolumn{2}{p{0.3\linewidth}}{
        \cellcolor[rgb]{0.9,0.9,0.9}{
            \makecell[{{p{\linewidth}}}]{
                \texttt{\tiny{[GM$|$GM]}}
                \texttt{calculator} \\
            }
        }
    }
    & & \\ \\

    \theutterance \stepcounter{utterance}  
    & & & \multicolumn{2}{p{0.3\linewidth}}{
        \cellcolor[rgb]{0.9,0.9,0.9}{
            \makecell[{{p{\linewidth}}}]{
                \texttt{\tiny{[GM$|$GM]}}
                ['şoarece', 'spaţiu']
            }
        }
    }
    & & \\ \\

    \theutterance \stepcounter{utterance}  
    & & & \multicolumn{2}{p{0.3\linewidth}}{
        \cellcolor[rgb]{0.9,0.9,0.9}{
            \makecell[{{p{\linewidth}}}]{
                \texttt{\tiny{[GM$|$GM]}}
                {'word': 'şoarece', 'assignment': 'team'}
            }
        }
    }
    & & \\ \\

    \theutterance \stepcounter{utterance}  
    & & & \multicolumn{2}{p{0.3\linewidth}}{
        \cellcolor[rgb]{0.9,0.9,0.9}{
            \makecell[{{p{\linewidth}}}]{
                \texttt{\tiny{[GM$|$GM]}}
                {'word': 'spaţiu', 'assignment': 'team'}
            }
        }
    }
    & & \\ \\

    \theutterance \stepcounter{utterance}  
    & & \multicolumn{4}{p{0.6\linewidth}}{
        \cellcolor[rgb]{0.9,0.9,0.9}{
            \makecell[{{p{\linewidth}}}]{
                \texttt{\tiny{[GM$\rangle$P2]}}
                \texttt{Furnizați o listă separată prin virgulă de până la 2 cuvinte din următoarea listă care se referă cel mai bine sau sunt asociate cel mai strâns cu cuvântul 'calculator'. Începeți întotdeauna lista de presupuneri cu GHICIRE:  și nu includeți niciun alt text în răspunsul dumneavoastră.} \\
\\ 
\texttt{pinguin, caracatiţă, pitic, germania, complot, bivol, canada, şoarece, împuşcat, prinţesă, rază, geniu, moarte, plastic, bloc, spaţiu, leu, mamut, cui, card, berlin, rădăcină, model, vânt, stadion} \\
            }
        }
    }
    & \\ \\

    \theutterance \stepcounter{utterance}  
    & & & \multicolumn{4}{p{0.6\linewidth}}{
        \cellcolor[rgb]{0.9,0.9,0.9}{
            \makecell[{{p{\linewidth}}}]{
                \texttt{\tiny{[GM$\langle$P2]}}
                \texttt{GHICIRE: şoarece, plastic} \\
            }
        }
    }
     \\ \\

    \theutterance \stepcounter{utterance}  
    & & & \multicolumn{2}{p{0.3\linewidth}}{
        \cellcolor[rgb]{0.9,0.9,0.9}{
            \makecell[{{p{\linewidth}}}]{
                \texttt{\tiny{[GM$|$GM]}}
                ['şoarece', 'plastic']
            }
        }
    }
    & & \\ \\

    \theutterance \stepcounter{utterance}  
    & & & \multicolumn{2}{p{0.3\linewidth}}{
        \cellcolor[rgb]{0.9,0.9,0.9}{
            \makecell[{{p{\linewidth}}}]{
                \texttt{\tiny{[GM$|$GM]}}
                {'word': 'şoarece', 'assignment': 'team'}
            }
        }
    }
    & & \\ \\

    \theutterance \stepcounter{utterance}  
    & & & \multicolumn{2}{p{0.3\linewidth}}{
        \cellcolor[rgb]{0.9,0.9,0.9}{
            \makecell[{{p{\linewidth}}}]{
                \texttt{\tiny{[GM$|$GM]}}
                {'word': 'şoarece', 'assignment': 'team'}
            }
        }
    }
    & & \\ \\

    \theutterance \stepcounter{utterance}  
    & & & \multicolumn{2}{p{0.3\linewidth}}{
        \cellcolor[rgb]{0.9,0.9,0.9}{
            \makecell[{{p{\linewidth}}}]{
                \texttt{\tiny{[GM$|$GM]}}
                {'word': 'plastic', 'assignment': 'team'}
            }
        }
    }
    & & \\ \\

    \theutterance \stepcounter{utterance}  
    & & & \multicolumn{2}{p{0.3\linewidth}}{
        \cellcolor[rgb]{0.9,0.9,0.9}{
            \makecell[{{p{\linewidth}}}]{
                \texttt{\tiny{[GM$|$GM]}}
                {'word': 'mamut', 'assignment': 'opponent'}
            }
        }
    }
    & & \\ \\

    \theutterance \stepcounter{utterance}  
    & & & \multicolumn{2}{p{0.3\linewidth}}{
        \cellcolor[rgb]{0.9,0.9,0.9}{
            \makecell[{{p{\linewidth}}}]{
                \texttt{\tiny{[GM$|$GM]}}
                {'hidden': {'team': ['leu', 'geniu', 'complot', 'spaţiu', 'rădăcină', 'canada', 'cui'], 'innocent': ['bloc', 'moarte', 'pitic'], 'opponent': ['caracatiţă', 'germania', 'vânt', 'stadion', 'prinţesă', 'pinguin', 'card'], 'assassin': ['bivol', 'împuşcat', 'rază', 'berlin', 'model']}, 'revealed': {'team': {'team': ['şoarece', 'plastic'], 'innocent': [], 'opponent': [], 'assassin': []}, 'opponent': {'team': [], 'innocent': [], 'opponent': ['mamut'], 'assassin': []}}}
            }
        }
    }
    & & \\ \\

    \theutterance \stepcounter{utterance}  
    & & \multicolumn{4}{p{0.6\linewidth}}{
        \cellcolor[rgb]{0.9,0.9,0.9}{
            \makecell[{{p{\linewidth}}}]{
                \texttt{\tiny{[P1$\langle$GM]}}
                \texttt{Cuvintele şoarece, plastic au fost ghicite corect. Rândul colegului tău de echipă s{-}a încheiat acolo.} \\
\texttt{Acum furnizați un alt indiciu referitor la unele dintre cuvintele echipei dvs. rămase și o listă a cuvintelor țintă aferente. Nu uitați să începeți indiciul cu 'INDICIU: ', puneți un rând nou și începeți lista de cuvinte țintă separate prin virgulă cu 'ȚINTE: '. Observație: unele cuvinte au fost eliminate din liste față de solicitările anterioare.} \\
\\ 
\texttt{Cuvintele echipei dvs. rămase sunt: leu, geniu, complot, spaţiu, rădăcină, canada, cui.} \\
\texttt{Cuvintele rămase pentru adversarul dvs. sunt: caracatiţă, germania, vânt, stadion, prinţesă, pinguin, mamut, card.} \\
\texttt{Cuvintele de distragere a atenției rămase sunt: bloc, moarte, pitic.} \\
\texttt{Cuvintele de ucidere rămase sunt: bivol, împuşcat, rază, berlin, model.} \\
            }
        }
    }
    & \\ \\

    \theutterance \stepcounter{utterance}  
    & \multicolumn{4}{p{0.6\linewidth}}{
        \cellcolor[rgb]{0.9,0.9,0.9}{
            \makecell[{{p{\linewidth}}}]{
                \texttt{\tiny{[P1$\rangle$GM]}}
                \texttt{INDICIU: uneltire} \\
\texttt{ȚINTE: complot, geniu} \\
            }
        }
    }
    & & \\ \\

    \theutterance \stepcounter{utterance}  
    & & & \multicolumn{2}{p{0.3\linewidth}}{
        \cellcolor[rgb]{0.9,0.9,0.9}{
            \makecell[{{p{\linewidth}}}]{
                \texttt{\tiny{[GM$|$GM]}}
                \texttt{uneltire} \\
            }
        }
    }
    & & \\ \\

    \theutterance \stepcounter{utterance}  
    & & & \multicolumn{2}{p{0.3\linewidth}}{
        \cellcolor[rgb]{0.9,0.9,0.9}{
            \makecell[{{p{\linewidth}}}]{
                \texttt{\tiny{[GM$|$GM]}}
                ['complot', 'geniu']
            }
        }
    }
    & & \\ \\

    \theutterance \stepcounter{utterance}  
    & & & \multicolumn{2}{p{0.3\linewidth}}{
        \cellcolor[rgb]{0.9,0.9,0.9}{
            \makecell[{{p{\linewidth}}}]{
                \texttt{\tiny{[GM$|$GM]}}
                {'word': 'complot', 'assignment': 'team'}
            }
        }
    }
    & & \\ \\

    \theutterance \stepcounter{utterance}  
    & & & \multicolumn{2}{p{0.3\linewidth}}{
        \cellcolor[rgb]{0.9,0.9,0.9}{
            \makecell[{{p{\linewidth}}}]{
                \texttt{\tiny{[GM$|$GM]}}
                {'word': 'geniu', 'assignment': 'team'}
            }
        }
    }
    & & \\ \\

    \theutterance \stepcounter{utterance}  
    & & \multicolumn{4}{p{0.6\linewidth}}{
        \cellcolor[rgb]{0.9,0.9,0.9}{
            \makecell[{{p{\linewidth}}}]{
                \texttt{\tiny{[GM$\rangle$P2]}}
                \texttt{Cuvintele şoarece, plastic au fost ghicite corect. Rândul tău s{-}a încheiat acolo.} \\
\texttt{Acum furnizați o altă listă separată prin virgulă de cel puțin 1 și până la 2 cuvinte din următoarea listă de cuvinte care se referă cel mai bine sau sunt cel mai strâns asociate cu cuvântul 'uneltire'. Nu uitați să vă începeți răspunsul cu GHICIRE: . Notă: unele cuvinte au fost eliminate din listă față de solicitările anterioare.} \\
\\ 
\texttt{pinguin, caracatiţă, pitic, germania, complot, bivol, canada, împuşcat, prinţesă, rază, geniu, moarte, bloc, spaţiu, leu, cui, card, berlin, rădăcină, model, vânt, stadion} \\
            }
        }
    }
    & \\ \\

    \theutterance \stepcounter{utterance}  
    & & & \multicolumn{4}{p{0.6\linewidth}}{
        \cellcolor[rgb]{0.9,0.9,0.9}{
            \makecell[{{p{\linewidth}}}]{
                \texttt{\tiny{[GM$\langle$P2]}}
                \texttt{GHICIRE: complot, geniu} \\
            }
        }
    }
     \\ \\

    \theutterance \stepcounter{utterance}  
    & & & \multicolumn{2}{p{0.3\linewidth}}{
        \cellcolor[rgb]{0.9,0.9,0.9}{
            \makecell[{{p{\linewidth}}}]{
                \texttt{\tiny{[GM$|$GM]}}
                ['complot', 'geniu']
            }
        }
    }
    & & \\ \\

    \theutterance \stepcounter{utterance}  
    & & & \multicolumn{2}{p{0.3\linewidth}}{
        \cellcolor[rgb]{0.9,0.9,0.9}{
            \makecell[{{p{\linewidth}}}]{
                \texttt{\tiny{[GM$|$GM]}}
                {'word': 'complot', 'assignment': 'team'}
            }
        }
    }
    & & \\ \\

    \theutterance \stepcounter{utterance}  
    & & & \multicolumn{2}{p{0.3\linewidth}}{
        \cellcolor[rgb]{0.9,0.9,0.9}{
            \makecell[{{p{\linewidth}}}]{
                \texttt{\tiny{[GM$|$GM]}}
                {'word': 'complot', 'assignment': 'team'}
            }
        }
    }
    & & \\ \\

    \theutterance \stepcounter{utterance}  
    & & & \multicolumn{2}{p{0.3\linewidth}}{
        \cellcolor[rgb]{0.9,0.9,0.9}{
            \makecell[{{p{\linewidth}}}]{
                \texttt{\tiny{[GM$|$GM]}}
                {'word': 'geniu', 'assignment': 'team'}
            }
        }
    }
    & & \\ \\

    \theutterance \stepcounter{utterance}  
    & & & \multicolumn{2}{p{0.3\linewidth}}{
        \cellcolor[rgb]{0.9,0.9,0.9}{
            \makecell[{{p{\linewidth}}}]{
                \texttt{\tiny{[GM$|$GM]}}
                {'word': 'geniu', 'assignment': 'team'}
            }
        }
    }
    & & \\ \\

    \theutterance \stepcounter{utterance}  
    & & & \multicolumn{2}{p{0.3\linewidth}}{
        \cellcolor[rgb]{0.9,0.9,0.9}{
            \makecell[{{p{\linewidth}}}]{
                \texttt{\tiny{[GM$|$GM]}}
                {'word': 'vânt', 'assignment': 'opponent'}
            }
        }
    }
    & & \\ \\

    \theutterance \stepcounter{utterance}  
    & & & \multicolumn{2}{p{0.3\linewidth}}{
        \cellcolor[rgb]{0.9,0.9,0.9}{
            \makecell[{{p{\linewidth}}}]{
                \texttt{\tiny{[GM$|$GM]}}
                {'hidden': {'team': ['leu', 'spaţiu', 'rădăcină', 'canada', 'cui'], 'innocent': ['bloc', 'moarte', 'pitic'], 'opponent': ['caracatiţă', 'germania', 'stadion', 'prinţesă', 'pinguin', 'card'], 'assassin': ['bivol', 'împuşcat', 'rază', 'berlin', 'model']}, 'revealed': {'team': {'team': ['şoarece', 'plastic', 'complot', 'geniu'], 'innocent': [], 'opponent': [], 'assassin': []}, 'opponent': {'team': [], 'innocent': [], 'opponent': ['mamut', 'vânt'], 'assassin': []}}}
            }
        }
    }
    & & \\ \\

    \theutterance \stepcounter{utterance}  
    & & \multicolumn{4}{p{0.6\linewidth}}{
        \cellcolor[rgb]{0.9,0.9,0.9}{
            \makecell[{{p{\linewidth}}}]{
                \texttt{\tiny{[P1$\langle$GM]}}
                \texttt{Cuvintele complot, geniu au fost ghicite corect. Rândul colegului tău de echipă s{-}a încheiat acolo.} \\
\texttt{Acum furnizați un alt indiciu referitor la unele dintre cuvintele echipei dvs. rămase și o listă a cuvintelor țintă aferente. Nu uitați să începeți indiciul cu 'INDICIU: ', puneți un rând nou și începeți lista de cuvinte țintă separate prin virgulă cu 'ȚINTE: '. Observație: unele cuvinte au fost eliminate din liste față de solicitările anterioare.} \\
\\ 
\texttt{Cuvintele echipei dvs. rămase sunt: leu, spaţiu, rădăcină, canada, cui.} \\
\texttt{Cuvintele rămase pentru adversarul dvs. sunt: caracatiţă, germania, vânt, stadion, prinţesă, pinguin, card.} \\
\texttt{Cuvintele de distragere a atenției rămase sunt: bloc, moarte, pitic.} \\
\texttt{Cuvintele de ucidere rămase sunt: bivol, împuşcat, rază, berlin, model.} \\
            }
        }
    }
    & \\ \\

    \theutterance \stepcounter{utterance}  
    & \multicolumn{4}{p{0.6\linewidth}}{
        \cellcolor[rgb]{0.9,0.9,0.9}{
            \makecell[{{p{\linewidth}}}]{
                \texttt{\tiny{[P1$\rangle$GM]}}
                \texttt{INDICIU: copac} \\
\texttt{ȚINTE: rădăcină, canada} \\
            }
        }
    }
    & & \\ \\

    \theutterance \stepcounter{utterance}  
    & & & \multicolumn{2}{p{0.3\linewidth}}{
        \cellcolor[rgb]{0.9,0.9,0.9}{
            \makecell[{{p{\linewidth}}}]{
                \texttt{\tiny{[GM$|$GM]}}
                \texttt{copac} \\
            }
        }
    }
    & & \\ \\

    \theutterance \stepcounter{utterance}  
    & & & \multicolumn{2}{p{0.3\linewidth}}{
        \cellcolor[rgb]{0.9,0.9,0.9}{
            \makecell[{{p{\linewidth}}}]{
                \texttt{\tiny{[GM$|$GM]}}
                ['rădăcină', 'canada']
            }
        }
    }
    & & \\ \\

    \theutterance \stepcounter{utterance}  
    & & & \multicolumn{2}{p{0.3\linewidth}}{
        \cellcolor[rgb]{0.9,0.9,0.9}{
            \makecell[{{p{\linewidth}}}]{
                \texttt{\tiny{[GM$|$GM]}}
                {'word': 'rădăcină', 'assignment': 'team'}
            }
        }
    }
    & & \\ \\

    \theutterance \stepcounter{utterance}  
    & & & \multicolumn{2}{p{0.3\linewidth}}{
        \cellcolor[rgb]{0.9,0.9,0.9}{
            \makecell[{{p{\linewidth}}}]{
                \texttt{\tiny{[GM$|$GM]}}
                {'word': 'canada', 'assignment': 'team'}
            }
        }
    }
    & & \\ \\

    \theutterance \stepcounter{utterance}  
    & & \multicolumn{4}{p{0.6\linewidth}}{
        \cellcolor[rgb]{0.9,0.9,0.9}{
            \makecell[{{p{\linewidth}}}]{
                \texttt{\tiny{[GM$\rangle$P2]}}
                \texttt{Cuvintele complot, geniu au fost ghicite corect. Rândul tău s{-}a încheiat acolo.} \\
\texttt{Acum furnizați o altă listă separată prin virgulă de cel puțin 1 și până la 2 cuvinte din următoarea listă de cuvinte care se referă cel mai bine sau sunt cel mai strâns asociate cu cuvântul 'copac'. Nu uitați să vă începeți răspunsul cu GHICIRE: . Notă: unele cuvinte au fost eliminate din listă față de solicitările anterioare.} \\
\\ 
\texttt{pinguin, caracatiţă, pitic, germania, bivol, canada, împuşcat, prinţesă, rază, moarte, bloc, spaţiu, leu, cui, card, berlin, rădăcină, model, stadion} \\
            }
        }
    }
    & \\ \\

    \theutterance \stepcounter{utterance}  
    & & & \multicolumn{4}{p{0.6\linewidth}}{
        \cellcolor[rgb]{0.9,0.9,0.9}{
            \makecell[{{p{\linewidth}}}]{
                \texttt{\tiny{[GM$\langle$P2]}}
                \texttt{GHICIRE: rădăcină, canada} \\
            }
        }
    }
     \\ \\

    \theutterance \stepcounter{utterance}  
    & & & \multicolumn{2}{p{0.3\linewidth}}{
        \cellcolor[rgb]{0.9,0.9,0.9}{
            \makecell[{{p{\linewidth}}}]{
                \texttt{\tiny{[GM$|$GM]}}
                ['rădăcină', 'canada']
            }
        }
    }
    & & \\ \\

    \theutterance \stepcounter{utterance}  
    & & & \multicolumn{2}{p{0.3\linewidth}}{
        \cellcolor[rgb]{0.9,0.9,0.9}{
            \makecell[{{p{\linewidth}}}]{
                \texttt{\tiny{[GM$|$GM]}}
                {'word': 'rădăcină', 'assignment': 'team'}
            }
        }
    }
    & & \\ \\

    \theutterance \stepcounter{utterance}  
    & & & \multicolumn{2}{p{0.3\linewidth}}{
        \cellcolor[rgb]{0.9,0.9,0.9}{
            \makecell[{{p{\linewidth}}}]{
                \texttt{\tiny{[GM$|$GM]}}
                {'word': 'rădăcină', 'assignment': 'team'}
            }
        }
    }
    & & \\ \\

    \theutterance \stepcounter{utterance}  
    & & & \multicolumn{2}{p{0.3\linewidth}}{
        \cellcolor[rgb]{0.9,0.9,0.9}{
            \makecell[{{p{\linewidth}}}]{
                \texttt{\tiny{[GM$|$GM]}}
                {'word': 'canada', 'assignment': 'team'}
            }
        }
    }
    & & \\ \\

    \theutterance \stepcounter{utterance}  
    & & & \multicolumn{2}{p{0.3\linewidth}}{
        \cellcolor[rgb]{0.9,0.9,0.9}{
            \makecell[{{p{\linewidth}}}]{
                \texttt{\tiny{[GM$|$GM]}}
                {'word': 'canada', 'assignment': 'team'}
            }
        }
    }
    & & \\ \\

    \theutterance \stepcounter{utterance}  
    & & & \multicolumn{2}{p{0.3\linewidth}}{
        \cellcolor[rgb]{0.9,0.9,0.9}{
            \makecell[{{p{\linewidth}}}]{
                \texttt{\tiny{[GM$|$GM]}}
                {'word': 'stadion', 'assignment': 'opponent'}
            }
        }
    }
    & & \\ \\

    \theutterance \stepcounter{utterance}  
    & & & \multicolumn{2}{p{0.3\linewidth}}{
        \cellcolor[rgb]{0.9,0.9,0.9}{
            \makecell[{{p{\linewidth}}}]{
                \texttt{\tiny{[GM$|$GM]}}
                {'hidden': {'team': ['leu', 'spaţiu', 'cui'], 'innocent': ['bloc', 'moarte', 'pitic'], 'opponent': ['caracatiţă', 'germania', 'prinţesă', 'pinguin', 'card'], 'assassin': ['bivol', 'împuşcat', 'rază', 'berlin', 'model']}, 'revealed': {'team': {'team': ['şoarece', 'plastic', 'complot', 'geniu', 'rădăcină', 'canada'], 'innocent': [], 'opponent': [], 'assassin': []}, 'opponent': {'team': [], 'innocent': [], 'opponent': ['mamut', 'vânt', 'stadion'], 'assassin': []}}}
            }
        }
    }
    & & \\ \\

    \theutterance \stepcounter{utterance}  
    & & \multicolumn{4}{p{0.6\linewidth}}{
        \cellcolor[rgb]{0.9,0.9,0.9}{
            \makecell[{{p{\linewidth}}}]{
                \texttt{\tiny{[P1$\langle$GM]}}
                \texttt{Cuvintele rădăcină, canada au fost ghicite corect. Rândul colegului tău de echipă s{-}a încheiat acolo.} \\
\texttt{Acum furnizați un alt indiciu referitor la unele dintre cuvintele echipei dvs. rămase și o listă a cuvintelor țintă aferente. Nu uitați să începeți indiciul cu 'INDICIU: ', puneți un rând nou și începeți lista de cuvinte țintă separate prin virgulă cu 'ȚINTE: '. Observație: unele cuvinte au fost eliminate din liste față de solicitările anterioare.} \\
\\ 
\texttt{Cuvintele echipei dvs. rămase sunt: leu, spaţiu, cui.} \\
\texttt{Cuvintele rămase pentru adversarul dvs. sunt: caracatiţă, germania, stadion, prinţesă, pinguin, card.} \\
\texttt{Cuvintele de distragere a atenției rămase sunt: bloc, moarte, pitic.} \\
\texttt{Cuvintele de ucidere rămase sunt: bivol, împuşcat, rază, berlin, model.} \\
            }
        }
    }
    & \\ \\

    \theutterance \stepcounter{utterance}  
    & \multicolumn{4}{p{0.6\linewidth}}{
        \cellcolor[rgb]{0.9,0.9,0.9}{
            \makecell[{{p{\linewidth}}}]{
                \texttt{\tiny{[P1$\rangle$GM]}}
                \texttt{INDICIU: constelaţie} \\
\texttt{ȚINTE: leu, spaţiu} \\
            }
        }
    }
    & & \\ \\

    \theutterance \stepcounter{utterance}  
    & & & \multicolumn{2}{p{0.3\linewidth}}{
        \cellcolor[rgb]{0.9,0.9,0.9}{
            \makecell[{{p{\linewidth}}}]{
                \texttt{\tiny{[GM$|$GM]}}
                \texttt{constelaţie} \\
            }
        }
    }
    & & \\ \\

    \theutterance \stepcounter{utterance}  
    & & & \multicolumn{2}{p{0.3\linewidth}}{
        \cellcolor[rgb]{0.9,0.9,0.9}{
            \makecell[{{p{\linewidth}}}]{
                \texttt{\tiny{[GM$|$GM]}}
                ['leu', 'spaţiu']
            }
        }
    }
    & & \\ \\

    \theutterance \stepcounter{utterance}  
    & & & \multicolumn{2}{p{0.3\linewidth}}{
        \cellcolor[rgb]{0.9,0.9,0.9}{
            \makecell[{{p{\linewidth}}}]{
                \texttt{\tiny{[GM$|$GM]}}
                {'word': 'leu', 'assignment': 'team'}
            }
        }
    }
    & & \\ \\

    \theutterance \stepcounter{utterance}  
    & & & \multicolumn{2}{p{0.3\linewidth}}{
        \cellcolor[rgb]{0.9,0.9,0.9}{
            \makecell[{{p{\linewidth}}}]{
                \texttt{\tiny{[GM$|$GM]}}
                {'word': 'spaţiu', 'assignment': 'team'}
            }
        }
    }
    & & \\ \\

    \theutterance \stepcounter{utterance}  
    & & \multicolumn{4}{p{0.6\linewidth}}{
        \cellcolor[rgb]{0.9,0.9,0.9}{
            \makecell[{{p{\linewidth}}}]{
                \texttt{\tiny{[GM$\rangle$P2]}}
                \texttt{Cuvintele rădăcină, canada au fost ghicite corect. Rândul tău s{-}a încheiat acolo.} \\
\texttt{Acum furnizați o altă listă separată prin virgulă de cel puțin 1 și până la 2 cuvinte din următoarea listă de cuvinte care se referă cel mai bine sau sunt cel mai strâns asociate cu cuvântul 'constelaţie'. Nu uitați să vă începeți răspunsul cu GHICIRE: . Notă: unele cuvinte au fost eliminate din listă față de solicitările anterioare.} \\
\\ 
\texttt{pinguin, caracatiţă, pitic, germania, bivol, împuşcat, prinţesă, rază, moarte, bloc, spaţiu, leu, cui, card, berlin, model} \\
            }
        }
    }
    & \\ \\

    \theutterance \stepcounter{utterance}  
    & & & \multicolumn{4}{p{0.6\linewidth}}{
        \cellcolor[rgb]{0.9,0.9,0.9}{
            \makecell[{{p{\linewidth}}}]{
                \texttt{\tiny{[GM$\langle$P2]}}
                \texttt{GHICIRE: spaţiu, leu} \\
            }
        }
    }
     \\ \\

    \theutterance \stepcounter{utterance}  
    & & & \multicolumn{2}{p{0.3\linewidth}}{
        \cellcolor[rgb]{0.9,0.9,0.9}{
            \makecell[{{p{\linewidth}}}]{
                \texttt{\tiny{[GM$|$GM]}}
                ['spaţiu', 'leu']
            }
        }
    }
    & & \\ \\

    \theutterance \stepcounter{utterance}  
    & & & \multicolumn{2}{p{0.3\linewidth}}{
        \cellcolor[rgb]{0.9,0.9,0.9}{
            \makecell[{{p{\linewidth}}}]{
                \texttt{\tiny{[GM$|$GM]}}
                {'word': 'spaţiu', 'assignment': 'team'}
            }
        }
    }
    & & \\ \\

    \theutterance \stepcounter{utterance}  
    & & & \multicolumn{2}{p{0.3\linewidth}}{
        \cellcolor[rgb]{0.9,0.9,0.9}{
            \makecell[{{p{\linewidth}}}]{
                \texttt{\tiny{[GM$|$GM]}}
                {'word': 'spaţiu', 'assignment': 'team'}
            }
        }
    }
    & & \\ \\

    \theutterance \stepcounter{utterance}  
    & & & \multicolumn{2}{p{0.3\linewidth}}{
        \cellcolor[rgb]{0.9,0.9,0.9}{
            \makecell[{{p{\linewidth}}}]{
                \texttt{\tiny{[GM$|$GM]}}
                {'word': 'leu', 'assignment': 'team'}
            }
        }
    }
    & & \\ \\

    \theutterance \stepcounter{utterance}  
    & & & \multicolumn{2}{p{0.3\linewidth}}{
        \cellcolor[rgb]{0.9,0.9,0.9}{
            \makecell[{{p{\linewidth}}}]{
                \texttt{\tiny{[GM$|$GM]}}
                {'word': 'leu', 'assignment': 'team'}
            }
        }
    }
    & & \\ \\

    \theutterance \stepcounter{utterance}  
    & & & \multicolumn{2}{p{0.3\linewidth}}{
        \cellcolor[rgb]{0.9,0.9,0.9}{
            \makecell[{{p{\linewidth}}}]{
                \texttt{\tiny{[GM$|$GM]}}
                {'word': 'germania', 'assignment': 'opponent'}
            }
        }
    }
    & & \\ \\

    \theutterance \stepcounter{utterance}  
    & & & \multicolumn{2}{p{0.3\linewidth}}{
        \cellcolor[rgb]{0.9,0.9,0.9}{
            \makecell[{{p{\linewidth}}}]{
                \texttt{\tiny{[GM$|$GM]}}
                {'hidden': {'team': ['cui'], 'innocent': ['bloc', 'moarte', 'pitic'], 'opponent': ['caracatiţă', 'prinţesă', 'pinguin', 'card'], 'assassin': ['bivol', 'împuşcat', 'rază', 'berlin', 'model']}, 'revealed': {'team': {'team': ['şoarece', 'plastic', 'complot', 'geniu', 'rădăcină', 'canada', 'spaţiu', 'leu'], 'innocent': [], 'opponent': [], 'assassin': []}, 'opponent': {'team': [], 'innocent': [], 'opponent': ['mamut', 'vânt', 'stadion', 'germania'], 'assassin': []}}}
            }
        }
    }
    & & \\ \\

    \theutterance \stepcounter{utterance}  
    & & \multicolumn{4}{p{0.6\linewidth}}{
        \cellcolor[rgb]{0.9,0.9,0.9}{
            \makecell[{{p{\linewidth}}}]{
                \texttt{\tiny{[P1$\langle$GM]}}
                \texttt{Cuvintele spaţiu, leu au fost ghicite corect. Rândul colegului tău de echipă s{-}a încheiat acolo.} \\
\texttt{Acum furnizați un alt indiciu referitor la unele dintre cuvintele echipei dvs. rămase și o listă a cuvintelor țintă aferente. Nu uitați să începeți indiciul cu 'INDICIU: ', puneți un rând nou și începeți lista de cuvinte țintă separate prin virgulă cu 'ȚINTE: '. Observație: unele cuvinte au fost eliminate din liste față de solicitările anterioare.} \\
\\ 
\texttt{Cuvintele echipei dvs. rămase sunt: cui.} \\
\texttt{Cuvintele rămase pentru adversarul dvs. sunt: caracatiţă, germania, prinţesă, pinguin, card.} \\
\texttt{Cuvintele de distragere a atenției rămase sunt: bloc, moarte, pitic.} \\
\texttt{Cuvintele de ucidere rămase sunt: bivol, împuşcat, rază, berlin, model.} \\
            }
        }
    }
    & \\ \\

    \theutterance \stepcounter{utterance}  
    & \multicolumn{4}{p{0.6\linewidth}}{
        \cellcolor[rgb]{0.9,0.9,0.9}{
            \makecell[{{p{\linewidth}}}]{
                \texttt{\tiny{[P1$\rangle$GM]}}
                \texttt{INDICIU: ciocan} \\
\texttt{ȚINTE: cui} \\
            }
        }
    }
    & & \\ \\

    \theutterance \stepcounter{utterance}  
    & & & \multicolumn{2}{p{0.3\linewidth}}{
        \cellcolor[rgb]{0.9,0.9,0.9}{
            \makecell[{{p{\linewidth}}}]{
                \texttt{\tiny{[GM$|$GM]}}
                \texttt{ciocan} \\
            }
        }
    }
    & & \\ \\

    \theutterance \stepcounter{utterance}  
    & & & \multicolumn{2}{p{0.3\linewidth}}{
        \cellcolor[rgb]{0.9,0.9,0.9}{
            \makecell[{{p{\linewidth}}}]{
                \texttt{\tiny{[GM$|$GM]}}
                ['cui']
            }
        }
    }
    & & \\ \\

    \theutterance \stepcounter{utterance}  
    & & & \multicolumn{2}{p{0.3\linewidth}}{
        \cellcolor[rgb]{0.9,0.9,0.9}{
            \makecell[{{p{\linewidth}}}]{
                \texttt{\tiny{[GM$|$GM]}}
                {'word': 'cui', 'assignment': 'team'}
            }
        }
    }
    & & \\ \\

    \theutterance \stepcounter{utterance}  
    & & \multicolumn{4}{p{0.6\linewidth}}{
        \cellcolor[rgb]{0.9,0.9,0.9}{
            \makecell[{{p{\linewidth}}}]{
                \texttt{\tiny{[GM$\rangle$P2]}}
                \texttt{Cuvintele spaţiu, leu au fost ghicite corect. Rândul tău s{-}a încheiat acolo.} \\
\texttt{Acum furnizați o altă listă separată prin virgulă de cel puțin 1 și până la 1 cuvinte din următoarea listă de cuvinte care se referă cel mai bine sau sunt cel mai strâns asociate cu cuvântul 'ciocan'. Nu uitați să vă începeți răspunsul cu GHICIRE: . Notă: unele cuvinte au fost eliminate din listă față de solicitările anterioare.} \\
\\ 
\texttt{pinguin, caracatiţă, pitic, bivol, împuşcat, prinţesă, rază, moarte, bloc, cui, card, berlin, model} \\
            }
        }
    }
    & \\ \\

    \theutterance \stepcounter{utterance}  
    & & & \multicolumn{4}{p{0.6\linewidth}}{
        \cellcolor[rgb]{0.9,0.9,0.9}{
            \makecell[{{p{\linewidth}}}]{
                \texttt{\tiny{[GM$\langle$P2]}}
                \texttt{GHICIRE: cui} \\
            }
        }
    }
     \\ \\

    \theutterance \stepcounter{utterance}  
    & & & \multicolumn{2}{p{0.3\linewidth}}{
        \cellcolor[rgb]{0.9,0.9,0.9}{
            \makecell[{{p{\linewidth}}}]{
                \texttt{\tiny{[GM$|$GM]}}
                ['cui']
            }
        }
    }
    & & \\ \\

    \theutterance \stepcounter{utterance}  
    & & & \multicolumn{2}{p{0.3\linewidth}}{
        \cellcolor[rgb]{0.9,0.9,0.9}{
            \makecell[{{p{\linewidth}}}]{
                \texttt{\tiny{[GM$|$GM]}}
                {'word': 'cui', 'assignment': 'team'}
            }
        }
    }
    & & \\ \\

    \theutterance \stepcounter{utterance}  
    & & & \multicolumn{2}{p{0.3\linewidth}}{
        \cellcolor[rgb]{0.9,0.9,0.9}{
            \makecell[{{p{\linewidth}}}]{
                \texttt{\tiny{[GM$|$GM]}}
                {'word': 'cui', 'assignment': 'team'}
            }
        }
    }
    & & \\ \\

    \theutterance \stepcounter{utterance}  
    & & & \multicolumn{2}{p{0.3\linewidth}}{
        \cellcolor[rgb]{0.9,0.9,0.9}{
            \makecell[{{p{\linewidth}}}]{
                \texttt{\tiny{[GM$|$GM]}}
                \texttt{team has won} \\
            }
        }
    }
    & & \\ \\

\end{supertabular}
}

\end{document}
